\documentclass{article}
\usepackage[utf8]{inputenc}
\usepackage[left=3cm, right=3cm, top=2cm, bottom=2cm]{geometry}
\usepackage[linesnumbered,boxed]{algorithm2e}
\usepackage[noend]{algpseudocode}
\usepackage{amsmath}
\usepackage{amssymb}
\usepackage{amsthm}
\usepackage{authblk}
\usepackage{comment}
\usepackage{mdframed}
\usepackage{graphicx}

\newtheorem{definition}{Definition}[section]

\begin{document}


\title{Homework 2 Solutions}
\author{CS 2051, Spring 2022}
\affil{Georgia Tech\\~\\{\bf Due Wed. Jan 26}}
\date{}
\maketitle

\paragraph{Reminder:} You may collaborate with other students in this
class, but (1) you must write up your own solutions in your own words, and (2) 
you must write down everyone you worked with at the top of the page, or ``no 
collaborators'' if you did it all on your own. Additionally, if you used any outside websites or textbooks besides the course text, please cite them here. You may not use question-answer sites like Chegg or MathOverflow. \\

You may assume the following definitions and anything we have defined in class:

\begin{definition}[Even numbers]
An integer $N$ is called even if there is an integer $K$ such that $2K = N$.
\end{definition}

\begin{definition}[Odd numbers]
An integer $N$ is called odd if there is an integer $K$ such that $2K + 1 = N$.
\end{definition}

\begin{definition}[Divisibility]
An integer $Z$ is said to divide $N$ if there exists an integer $K$ such that $ZK = N$.
\end{definition}


\hrulefill

\begin{enumerate}

    \item For each argument below, label each of the statements and write the premises and the conclusion using propositional variables and connectives. Then, check the validity of the arguments. You should properly label the rule of inference you use on each step along with the lines to which you apply it. 
    \begin{enumerate}
        \item "Chris, a Georgia Tech graduate, knows how to write proofs. Everyone who can write proofs can get a high-paying job. Therefore, there exists a person who graduates from Georgia Tech who can get a high paying job." 
        
        \begin{mdframed}
        Let:\\
        $P(x) = $ "x is a student of this class" \\
        $Q(x) = $ "x knows how to write proofs" \\
        $R(x) = $ "x can get a high-paying job" \\
        \\
    The argument reads:
\begin{align*}
& (1) \quad ~P(Doug)\wedge Q(Doug) \\
& (2) \quad ~\forall x(Q(x)\rightarrow R(x))\\
& \noindent\rule{4cm}{0.4pt}\\
& \therefore \quad ~ \exists x(P(x)\wedge R(x))
\end{align*}

We can check this argument as follows: \\
\begin{align*}
& (3) \quad ~Q(Doug) \mbox{ from $(1)$ using simplification.} \\
& (4) \quad ~Q(Doug)\rightarrow R(Doug) \mbox{ from $(2)$ using universal instantiation.}\\
& (5) \quad ~R(Doug) \mbox{ from $(3),~(4)$ using modus ponens.}\\
& (6) \quad ~P(Doug) \mbox{ from $(1)$ using simplification.}\\
& (7) \quad ~P(Doug)\wedge R(Doug) \mbox{ from $(5),~(6)$ using conjunction.}\\
& (8) \quad ~ \exists x(P(x)\wedge R(x)) \mbox{ from $(7)$ using existential generalization.}
\end{align*}
        \end{mdframed}
        \item "Someone in this class enjoys whale watching. Every person who enjoys whale watching cares about ocean pollution. Therefore, there is a person in this class who cares about ocean pollution."
        \begin{mdframed}
        
        {Let:\\
$P(x)$  = "x is a student of this class" \\
$Q(x)$ = "x enjoys whale watching" \\
$R(x)$ = "x cares about ocean pollution"\\
        
            The argument reads:

\begin{align*}
& (1) \quad ~\exists x(P(x)\wedge Q(x))) \\
& (2) \quad ~\forall x(Q(x)\rightarrow R(x))\\
& \noindent\rule{4cm}{0.4pt}\\
& \therefore \quad ~ \exists x (P(x)\wedge R(x))
\end{align*}

We can check this argument as follows:\\
\begin{align*}
& (3) \quad ~P(c)\wedge Q(c) \mbox{ from $(1)$ using existential instantiation.}\\
& (4) \quad ~Q(c) \mbox{ from $(3)$ using existential simplification.}\\
& (5) \quad ~Q(c)\rightarrow R(c) \mbox{ from $(2)$ using universal instantiation.}\\
& (6) \quad ~R(c) \mbox{ from $(4),~(5)$ using modus ponens.}\\
& (7) \quad ~P(c) \mbox{ from $(3)$ using simplification.}\\
& (8) \quad ~P(c)\wedge R(c) \mbox{ from $(6),~(7)$ using conjunction.}
\end{align*}
        \end{mdframed}
    \end{enumerate}
    
    \item Show that the following argument is valid:
    
\begin{align*}
& (1) \quad ~(p \wedge t) \rightarrow (r \vee s)\\
& (2) \quad ~q \rightarrow (u\wedge t)\\
& (3) \quad ~u \rightarrow  p\\
& (4) \quad ~\neg ~ s\\
&\noindent\rule{4cm}{0.4pt}\\
& \therefore \quad ~ q\rightarrow r
\end{align*}    


\begin{mdframed}
We can break the problem into cases where $q$ is true or $q$ is false:
\begin{enumerate}
    \item \textbf{q is false}. Then $q \rightarrow r = \neg q \lor r$ is true, so we are done.
    \item \textbf{q is true}. Then we can rewrite the original premise as:
    
\begin{align*}
& (1) \quad ~(p \wedge t) \rightarrow (r \vee s)\\
& (2) \quad ~q \rightarrow (u\wedge t)\\
& (3) \quad ~u \rightarrow  p\\
& (4) \quad ~\neg ~ s\\
& (5) \quad ~q\\
&\noindent\rule{4cm}{0.4pt}\\
& \therefore \quad ~ r
\end{align*}    

    And we check the argument as follows:

\begin{align*}
& (6) \quad ~(u\wedge t) \mbox{ from $(2), ~(5)$ using modus ponens.}\\
& (7) \quad ~u \mbox{ from $(6)$ using simplification.}\\
& (8) \quad ~p \mbox{ from $(3), ~(7)$ using modus ponens.}\\
& (9) \quad ~t \mbox{ from $(6)$ using simplification.}\\
& (10) \quad ~(p\wedge t) \mbox{ from $(8),~(9)$ using conjunction.}\\
& (11) \quad ~(r\vee s) \mbox{ from $(1),~(10)$ using modus ponens.}\\
& (12) \quad ~ r \mbox{ from $(4),~(11)$ using disjunctive syllogism.}\\
\end{align*}
\end{enumerate}
\end{mdframed}


    \item Prove the \textbf{triangle inequality}, which states that if $x, y \in \mathbb{R}$, then $|x| + |y| \geq |x + y|$. Here, $|x|$ is the absolute value of $x$ defined as
\[
|x|=\begin{cases}
& x, \mbox{ if $x\geq 0$}\\
& -x \mbox{ if $x< 0$}
\end{cases}
\]
\textit{Hint: break the proof up into cases, considering all possible choices of $x$ and $y$ being positive, negative, or zero.}
    
    \begin{mdframed}
        First, observe that $|x+y|=|y+x|$ since addition is commutative, thus we can assume that $x\geq y$ without a loss of generality.
    Now, proceed with cases:
    \begin{enumerate}
        \item Case 1: $x=0$. Then $|x+y|=|y|=0+|y|=|x|+|y|$.
        \item Case 2: $y=0$. This is analogous to the previous case.
        \item Case 3: $x\geq y>0$. Then, by definition of absolute value, $|x+y|=x+y=|x|+|y|$.
        \item Case 4: $x>0>y$. We further note that $|x+y|=|-x-y|$, thus, we can further assume that $|x|>|y|$. Then $|x+y|=x-|y|=|x|-|y|<|x|+|y|$.
        \item Case 5: $0>x>y$. Follows from case 3 and the observation $|x+y|=|-x-y|$.
    \end{enumerate}
    Thus, the triangle inequality holds true.
    \end{mdframed}
    
    \item Show that if $n$ is an integer and $3(n + 4)^2 + 2$ is even, then $n$ is even.
    \begin{enumerate}
        \item using a proof by contradiction
        \begin{mdframed}
            Assume for contradiction that if $3(n + 4)^2 + 2$ is even, then $n$ is odd. Then $n = 2j + 1$ for some $j \in \mathbb{Z}$. $3(n+4)^2+2 = 3(n^2 + 8n + 16) + 2 = 3n^2 + 24n + 18 = 3(2j + 1)^2 + 24(2j + 1) + 18 = 3(4j^2 + 4j + 1) + 48j + 24 + 18 = 12j^2 + 60j + 48 + 3 = 2(6j^2 + 30 + 24 + 1) + 1$. Now, $p = 6j^2 + 30 + 24 + 1 \in \mathbb{Z}$, so $2(6j^2 + 30 + 24 + 1) + 1 = 2p + 1$, so $3(n+4)^2 + 2$ is odd, which is a contradiction. Therefore, $n$ is even.
        \end{mdframed}
        \item using a proof by contrapositive
            \begin{mdframed}
            The contrapositive is: if $n$ is odd, then $3(n+4)^2 + 2$ is odd. \\
            $n = 2j + 1$ because it is odd. So $3(n+4)^2 + 2$ is odd by the same algebraic manipulation as in part (a).
            \end{mdframed}
    \end{enumerate}

    \item In this problem, you will prove that $\sqrt{3}$ is irrational.
    \begin{enumerate}
        \item Prove the following lemma: if $a \in \mathbb{Z}$ and $a = 3k + 1$ for some $k \in \mathbb{Z}$, then $a^2 = 3j + 1$ for some integer $j$.
        \begin{mdframed}
        Direct proof: $a^2 = (3k+1)^2 = 9k^2 + 6k + 1 = 3(3k^2 + 2k) + 1$. $3k^2 + 2k \in \mathbb{Z}$, so $a^2 = 3j + 1$ when $j = 3k^2 + 2k$
        \end{mdframed}
        \item Prove the following lemma: if $a \in \mathbb{Z}$ and $a = 3k + 2$ for some $k \in \mathbb{Z}$, then $a^2 = 3j + 1$ for some integer $j$.
        \begin{mdframed}
        Direct proof: $a^2 = (3k+2)^2 = 9k^2 + 12k + 4 = 9k^2 + 12k + 3 + 1 = 3(3k^2 + 4k + 1) + 1$. $3k^2 + 4k + 1\in \mathbb{Z}$, so $a^2 = 3j + 2$ when $j = 3k^2 + 4k + 1$.
        \end{mdframed}
        \item Prove the following lemma: if $a^2$ is an integer and $a^2$ is divisible by 3, then $a$ is also divisible by 3. (Hint: use the fact that exactly one of the following statements is true: there is an integer $k$ such that $a = 3k$, there is an integer $k$ such that $a = 3k + 1$, or there is an integer $k$ such that $a = 3k + 2$).
        \begin{mdframed}
        Prove the contrapositive: If $a$ is not divisible by 3, then $a^2$ is not divisible by 3. \\
        Now, break into the two cases of not being divisible by three proved above:
        \begin{enumerate}
            \item $a = 3k + 1$. Then by lemma (a), $a^2 = 3j + 1$ which is not divisible by 3.
            \item $a = 3k + 2$. Then by lemma (a), $a^2 = 3j + 1$ which is not divisible by 3.
        \end{enumerate}
        
        \end{mdframed}
        \item Use the previous lemma to prove that $\sqrt{3}$ is irrational. (Hint: use a proof by contradiction. Assume that $\sqrt{3}$ is rational, let $k$ be the smallest positive integer such that $k\sqrt{3}$ is an integer, then arrive at a contradiction by proving that some other integer $j < k$ also works). 
        \begin{mdframed}
        Assume towards contradiction that $\sqrt{3}$ is rational and let $k$ be the smallest positive integer such that $k\sqrt{3}$ is also an integer. Then, let $a = k\sqrt{3}$ so that $a^2 = 3k^2$. By lemma (c), because $a^2$ is divisible by 3, $a$ is also divisible by 3. Therefore, $\frac{a}{3} = \sqrt{3}\frac{k}{3}$ is an integer. Let $j = \frac{k}{3}$. $j < k$, so $k$ is not the smallest integer such that $k\sqrt{3}$ is also an integer, which is a contradiction.
        \end{mdframed}
        \item 4 is clearly rational because $\sqrt{4} = 2$. Explain which step goes wrong in your above proof on $\sqrt{4}$ instead of $\sqrt{3}$.
        \begin{mdframed}
        Lemma (c) does not hold. If $a^2$ is an integer and $a^2$ is divisible by 4, then $a$ is not necessarily divisible by 4. As a counterexample, $a^2 = 4$, then $a = 2$ which is not divisible by 4.
        \end{mdframed}
    \end{enumerate}
        
    
\end{enumerate}



\end{document}
